\chapter{Worked example: $\Sigma$-protocol zero knowledge}

The zero-knowledge side of Schnorr, complementing the forking soundness of
Chapter~11. For the abstract $\Sigma$-protocol (instantiated by the XOR relation
$\mathsf{simpleSigma}$), all four standard properties hold, and the commitment
hiding/binding pair follows. Each proof is a few lines: the coupling toolkit and
$\mathsf{ssprove\_distr\_simp}$ reduce the arguments to a single bijection couple
or a finite case split.

\begin{theorem}[Completeness]
  \label{thm:sigma-complete}
  \uses{def:spcomp}
  \lean{CatCrypt.Examples.SigmaProtocol.simpleSigma_complete}
  \leanok
  Honest execution with a valid witness always verifies.
\end{theorem}

\begin{theorem}[Special honest-verifier zero knowledge]
  \label{thm:sigma-shvzk}
  \uses{thm:sigma-complete}
  \lean{CatCrypt.Examples.SigmaProtocol.simpleSigma_shvzk, CatCrypt.Examples.SigmaProtocol.simpleSigma_shvzk_zero_adv}
  \leanok
  Real and simulated transcripts are identically distributed: the simulator, given
  the challenge, produces a transcript coupled to the real one through a single
  bijection, so the distinguishing advantage is $0$.
\end{theorem}

\begin{theorem}[Special soundness]
  \label{thm:sigma-sound}
  \uses{thm:sigma-complete}
  \lean{CatCrypt.Examples.SigmaProtocol.simpleSigma_special_sound}
  \leanok
  Two accepting transcripts sharing a commitment but differing in challenge yield
  the witness through an explicit extractor.
\end{theorem}

\begin{theorem}[Commitment hiding and binding]
  \label{thm:sigma-commitment}
  \uses{thm:sigma-shvzk, thm:sigma-sound}
  \lean{CatCrypt.Examples.SigmaProtocol.simpleSigma_hiding, CatCrypt.Examples.SigmaProtocol.simpleSigma_binding}
  \leanok
  The associated commitment is perfectly hiding (a direct instance of SHVZK's
  one-time-pad coupling) and binding (a double opening yields two accepting
  transcripts, hence the witness by special soundness).
\end{theorem}

\section{Stateful provers and the given-challenge formulation}

A second module develops the same three properties with a prover that carries
state from the first move to the third, together with conversions between the
stateful and the stateless presentations. SHVZK is stated both for a sampled
challenge and for a challenge fixed in advance.

\begin{definition}[$\Sigma$-protocol with prover state]
  \label{def:sigma-stateful}
  \uses{def:spcomp}
  \lean{CatCrypt.Examples.Sigma.SigmaProtocol, CatCrypt.Examples.Sigma.SigmaProtocolS, CatCrypt.Examples.Sigma.SHVZK_given, CatCrypt.Examples.Sigma.SpecialSoundness}
  \leanok
  A three-move protocol whose prover state passes from the commitment to the
  response, with the completeness, special-soundness and SHVZK predicates.
\end{definition}

\begin{theorem}[The XOR protocol, stateful presentation]
  \label{thm:sigma-stateful-xor}
  \uses{def:sigma-stateful}
  \lean{CatCrypt.Examples.Sigma.simpleSigma_shvzk_given, CatCrypt.Examples.Sigma.simpleSigma_shvzk, CatCrypt.Examples.Sigma.simpleSigma_special_sound}
  \leanok
  The XOR relation protocol is SHVZK for a given challenge and for a sampled
  challenge, and is specially sound.
\end{theorem}
